\begin{figure}[H]
\centering
\begin{tikzpicture}[
  node distance=0mm,
  cell/.style={rectangle, draw, minimum size=8mm, font=\ttfamily}
]
  % Pole
  \node[cell] (c1) {3};
  \node[cell, right=of c1] (c2) {5};
  \node[cell, right=of c2] (c3) {7};
  \node[cell, right=of c3] (c4) {11};
  \node[cell, right=of c4] (c5) {6};
  \node[cell, right=of c5] (c6) {9};
  
  % Indexy
  \node[above=1mm of c1, font=\scriptsize\sffamily] {1};
  \node[above=1mm of c2, font=\scriptsize\sffamily] {2};
  \node[above=1mm of c3, font=\scriptsize\sffamily] {3};
  \node[above=1mm of c4, font=\scriptsize\sffamily] {4};
  \node[above=1mm of c5, font=\scriptsize\sffamily] {5};
  \node[above=1mm of c6, font=\scriptsize\sffamily] {6};
\end{tikzpicture}
\caption{Reprezentace haldy v poli $H[1 \dots 6]$. Například uzel na indexu $i=2$ s hodnotou $5$ má levého syna na indexu $2\cdot 2 = 4$ (hodnota 11) a pravého syna na indexu $2\cdot 2 + 1 = 5$ (hodnota 6).}
\end{figure}